\section*{第 \ref{ch:b1:filters-transfer-functions} 章——滤波器与传递函数}

\begin{solution}{exo:b1:filters-transfer-functions:1}
$0.5 \to \qty{-6.0}{dB}$; $10 \to \qty{20}{dB}$; $1/\sqrt2 \to \qty{-3.0}{dB}$;
$100 \to \qty{40}{dB}$. $\qty{-40}{dB} \to 0.01$; $\qty{+6}{dB} \to 2.0$.
级联：$-6 - 20 = \qty{-26}{dB}$，比值为 $0.05$。
\end{solution}

\begin{solution}{exo:b1:filters-transfer-functions:2}
$f_c = 1/(2\pi RC) = 1/(2\pi \times 10^4 \times 1.6\times10^{-8}) =
\qty{1.0}{kHz}$。增益 $-10\log(1 + (f/f_c)^2)$：\qty{100}{Hz}：
\qty{-0.04}{dB}; \qty{1}{kHz}: \qty{-3}{dB}; \qty{10}{kHz}: \qty{-20}{dB};
\qty{100}{kHz}: \qty{-40}{dB}.
\end{solution}

\begin{solution}{exo:b1:filters-transfer-functions:3}
(a) 低通滤波器（$C$ 在 $f$ 较高时使输出短路）；(b) 高通滤波器（$C$
在 $f$ 较低时断开）；(c) 二阶低通滤波器（$L$ 在 $f$ 较高时断开，
$C$ 使其短路）；(d) 带通滤波器（$C$ 阻断 $f$ 较低的成分，$L$ 阻断 $f$ 较高的成分）。
\end{solution}

\begin{solution}{exo:b1:filters-transfer-functions:4}
$4E/n\pi$：在 \qty{1}{kHz} 时为 \qty{6.37}{V}，在 \qty{3}{kHz} 时为 \qty{2.12}{V}，
在 \qty{5}{kHz} 时为 \qty{1.27}{V}。方波在
半周期平移下呈反对称（$s(t + T/2) = -s(t)$），而偶次谐波
则不具有此性质。
\end{solution}

\begin{solution}{exo:b1:filters-transfer-functions:5}
在 \qty{100}{Hz} 及其有用谐波处，$\omega \ll \omega_c$：
$\underline H \approx j\omega/\omega_c$，即导数除以
$\omega_c$。三角波的斜率为 $\pm 4Ef = \pm\qty{400}{V/s}$，故
输出为方波，其振幅为 $400/(2\pi \times 10^4) =
\qty{6.4}{mV}$。
\end{solution}

\begin{solution}{exo:b1:filters-transfer-functions:6}
$\omega \gg \omega_c$: $u_s \approx (1/\tau)\int u_e\,\dd t$, $\tau =
1/(2\pi \times 10) = \qty{16}{ms}$。$\pm E$ 的积分是一个
三角波，每半周期上升 $ET/2$，摆幅为 $ET/2\tau$，振幅
为 $ET/4\tau = 1.0 \times 10^{-3}/(4 \times 0.016) = \qty{16}{mV}$，围绕
零值。
\end{solution}

\begin{solution}{exo:b1:filters-transfer-functions:7}
$f_0 = \qty{5.03}{kHz}$. $Q = \sqrt{L/C}/R = 1/\sqrt2$: $R = \sqrt2 \times
316 = \qty{447}{\ohm}$。在 $f_0$ 处：\qty{-3}{dB}；在 $10f_0$ 处：$1/\sqrt{1 +
10^4} \to \qty{-40}{dB}$.
\end{solution}

\begin{solution}{exo:b1:filters-transfer-functions:8}
$G = 1/\sqrt{1 + Q^2(x - 1/x)^2}$，其中 $x = f/3$；在 \qty{1}{kHz} 处：$x = 1/3$，
$Q(x - 1/x) = -53$, $G = 0.019$; \qty{3}{kHz}: $G = 1$; \qty{5}{kHz}:
$x = 5/3$，$21$，$G = 0.047$。输出在 \qty{3}{kHz} 处为 \qty{2.1}{V}，其中
在 \qty{1}{kHz} 处为 \qty{0.12}{V}，在 \qty{5}{kHz} 处为 \qty{0.06}{V}——基本上
是一个近似纯的 \qty{3}{kHz} 正弦波。
\end{solution}

\begin{solution}{exo:b1:filters-transfer-functions:9}
节点法：令 $a = jRC\omega$，第一个电容器在上游看到 $R$，
在下游看到 $R + 1/jC\omega$，于是 $\underline H = 1/(1 + 3a + a^2)$，即\
$1/[1 - (\omega/\omega_c)^2 + 3j\omega/\omega_c]$。在 $\omega_c$ 处：$1/3j$，
$G = 1/3$，\qty{-9.5}{dB}；两个独立增益之积应为
$1/2$, \qty{-6}{dB}. $f_c = \qty{1.0}{kHz}$.
\end{solution}

\begin{solution}{exo:b1:filters-transfer-functions:10}
低于 $\omega_1$ 时为 \qty{0}{dB}；介于 $\omega_1$ 与 $\omega_2$ 之间时，增益
按 \qty{20}{dB/decade} 上升；高于 $\omega_2$ 时保持平坦，为 $20\log(\omega_2/\omega_1) =
\qty{40}{dB}$。相位为 $\arctan(\omega/\omega_1) - \arctan(\omega/\omega_2)$：
在 $\omega = 10^3$ 处，$\arctan 10 - \arctan 0.1 = 84^\circ - 6^\circ =
78^\circ$.
\end{solution}

\begin{solution}{exo:b1:filters-transfer-functions:11}
远高于截止频率时，$G \approx f_c/f$；滤波后的纹波比
为
\[
\frac{(4E/3\pi)(f_c/100)}{2E/\pi} = \frac{2}{3}\,\frac{f_c}{100} < 0.01,
\]
因此 $f_c < \qty{1.5}{Hz}$，且 $\tau = 1/2\pi f_c > \qty{0.1}{s}$。负载本身即为
$R$：并联在其两端的电容不会造成电压降与功率损耗，
而串联电阻则两者都会被白白消耗。
\end{solution}

\begin{solution}{exo:b1:filters-transfer-functions:12}
$f_c = 1/(2\pi RC) = 1/(2\pi \times 1.6\times10^5 \times 10^{-7}) =
\qty{10}{Hz}$，$\tau = RC = \qty{16}{ms}$。平顶段内
电容器通过 $R$ 放电：输出按 $\eu^{-t/\tau}$ 衰减，
自每个边沿起。半周期 $T/2$ 内的跌落为 $1 - \eu^{-T/2\tau}$：\qty{50}{Hz}
($T/2 = \qty{10}{ms}$): $1 - \eu^{-0.63} = 47\%$; \qty{1}{kHz}
($T/2 = \qty{0.5}{ms}$): $3\%$.
\end{solution}

\begin{solution}{pb:b1:filters-transfer-functions:1}
\textbf{1.} $\underline H = \underline U_{\text{驱动}}/\underline U_e$；
$G_{\mathrm{dB}} = 20\log\abs{\underline H}$.

\textbf{2.} $\underline H_L = R/(R + jL\omega) = 1/(1 + j\omega/\omega_c)$,
$\omega_c = R/L$.

\textbf{3.} $L = R/\omega_c = 8.0/(2\pi \times 2000) = \qty{0.64}{mH}$.

\textbf{4.} $\underline H_H = R/(R + 1/jC\omega) = jRC\omega/(1 + jRC\omega)$,
$\omega_c = 1/RC$; $C = 1/(R\omega_c) = \qty{9.9}{\micro F} \approx
\qty{10}{\micro F}$.

\textbf{5.} 在 $f_c$ 处各为 \qty{-3}{dB}。\qty{200}{Hz}（$x = 0.1$）时：低音扬声器
\qty{-0.04}{dB}，高音扬声器 \qty{-20}{dB}。\qty{20}{kHz}（$x = 10$）时：低音扬声器
\qty{-20}{dB}，高音扬声器 \qty{-0.04}{dB}。

\textbf{6.} $G_L^2 + G_H^2 = (1 + x^2)/(1 + x^2) = 1$：负载相等时
两者功率之和等于单个 \qty{8}{\ohm} 驱动单元
所吸收的功率——是分配而非损失。

\textbf{7.} $-45^\circ$（低音单元）与 $+45^\circ$（高音单元）：相差 $90^\circ$；
在 $f_c$ 处两个纸盆不同步运动。

\textbf{8.} 每倍频程 \qty{\pm 6}{dB}（每十倍频程 \qty{20}{dB}）。

\textbf{9.} 由 $jL\omega$ 与 $R \parallel 1/jC\omega =
R/(1 + jRC\omega)$ 构成的分压：$\underline H = R/[R + jL\omega(1 + jRC\omega)] =
1/(1 - LC\omega^2 + jL\omega/R)$。

\textbf{10.} $\omega_0 = 1/\sqrt{LC}$；由 $L\omega/R = x/Q$ 得 $Q =
R/L\omega_0 = R\sqrt{C/L}$。

\textbf{11.} $\abs{\underline H}^{-2} = (1 - x^2)^2 + x^2/Q^2 = 1 + x^4 +
x^2(1/Q^2 - 2) = 1 + x^4$，当 $Q^2 = \tfrac12$ 时；$x = 1$ 时 $G = 1/\sqrt2$。

\textbf{12.} $L = R/(Q\omega_0) = 8\sqrt2/(1.257\times10^4) = \qty{0.90}{mH}$；
$C = 1/(L\omega_0^2) = \qty{7.0}{\micro F}$。

\textbf{13.} $x = 2$ 时：$1/\sqrt{17}$，\qty{-12.3}{dB}；$x = 4$ 时：$1/\sqrt{257}$，
\qty{-24.1}{dB}：斜率为每倍频程 \qty{12}{dB}。

\textbf{14.} 交换 $L$ 与 $C$ 即交换 $x \leftrightarrow 1/x$（而
$Q$ 不变）：$\underline H_H = -x^2/(1 - x^2 + jx/Q)$。

\textbf{15.} 在 $x = 1$ 处：$\underline H_L = Q/j = -jQ$（$-90^\circ$），
$\underline H_H = jQ$（$+90^\circ$）：两个单元在分频点处
反相，声学上会相互抵消；高音单元
应接成极性反接。

\textbf{16.} 高于 \qty{2}{kHz} 的谐音：$n \geq 10$（\qty{2.2}{kHz}
及以上）主要进入高音单元。

\textbf{17.} $9 \times 220 = \qty{1980}{Hz} \approx f_c$：平均分配，
两者各得 \qty{-3}{dB}。

\textbf{18.} 低音单元（$G_L$）：$x = 0.5$ 时 $0.89$；$x = 1.5$ 时 $0.55$；$x = 2.5$ 时
$0.37$。高音单元（$G_H$）：$0.45$、$0.83$、$0.93$。

\textbf{19.} 高音单元：其串联电容器隔断直流。低音
单元承受偏置（\qty{1}{V} 加在 \qty{8}{\ohm} 两端，电流 \qty{125}{mA}，
发热 \qty{0.125}{W}，纸盆偏移）。

\textbf{20.} 一阶，$x = 0.25$：$G_H = 0.25/\sqrt{1.0625} = 0.24$
（\qty{-12}{dB}），功率占比 $6\%$。二阶：$0.0625/\sqrt{1 +
0.0039} = 0.06$（\qty{-24}{dB}），$0.4\%$。额定功率仅几瓦的高音单元
若采用一阶设计，会接收到一个强低音音符的 $6\%$——它可能
被烧毁。

\textbf{21.} $\underline Z = 8 + jL_{vc}\omega$：\qty{2}{kHz} 时 $8 + 6.3j$，
$\abs{\underline Z} = \qty{10.2}{\ohm}$；\qty{8}{kHz} 时 $8 + 25.1j$，
\qty{26.3}{\ohm}。

\textbf{22.} $\underline H = \underline Z/(\underline Z + jL\omega)$：
\qty{2}{kHz} 时：$10.2/\abs{8 + 14.3j} = 10.2/16.4 = 0.62$（\qty{-4.2}{dB}）；
\qty{8}{kHz} 时：$26.3/\abs{8 + 57.3j} = 0.45$（\qty{-6.9}{dB}），而非
\qty{-12.3}{dB}：升高的线圈阻抗与串联线圈相抗衡，
滚降变缓，每倍频程仅几分贝。

\textbf{23.} 导纳：$1/(R + jL_{vc}\omega) + jC_z\omega/(1 + jRC_z\omega)$；
取 $C_z = L_{vc}/R^2$，则第二项为 $jL_{vc}\omega/[R(R +
jL_{vc}\omega)]$，二者之和为 $(R + jL_{vc}\omega)/[R(R + jL_{vc}\omega)]
= 1/R$。$C_z = 0.50\times10^{-3}/64 = \qty{7.8}{\micro F}$。

\textbf{24.} $U = \sqrt{PR} = \sqrt{400} = \qty{20}{V}$ 有效值。在 $f_c$ 处每个
单元各得一半：\qty{25}{W}。在 \qty{500}{Hz} 处高音单元获得 $6\%$：
\qty{3}{W}。

\textbf{25.} 一阶：两部分，功率互补，平缓的
\qty{6}{dB/oct} 斜率，单元之间相差 $90^\circ$，低音泄漏到
高音单元。二阶：四个部分，\qty{12}{dB/oct}，分频点处相差 $180^\circ$
（需将高音单元反接），对高音单元的保护好得多。两种
设计中，低音单元的电感都必须用 佐贝尔网络加以抑制，
方案才能成立。
\end{solution}
